%% ch04 --- Ruch: reguły zmiany i mapa wszystkich stanów
%%
%% Napisany we wrześniu 2026. Każda liczba, której czytelnik nie policzy w
%% pamięci, pochodzi z code/measure/ch04.py, który importuje TE SAME funkcje,
%% które drukują listingi (code/ch04/ i flows z chaoslab), więc strona i
%% listing nie mogą się rozejść. Sprężyna, pętla i planeta używają tylko
%% + - * / i sqrt, więc ich cyfry są takie same na każdej maszynie; wahadło
%% woła math.sin, ale nie jest chaotyczne, a jego czasy wahnięć są drukowane
%% z trzema miejscami po przecinku, których różnica w ostatnim bicie nie
%% ruszy.
%%
%% Bliźniak: chapters/en/ch04-motion.tex. Podrozdział za podrozdziałem, makro
%% za makrem, liczba za liczbą; tools/parity.py je porównuje.

\chapter{Ruch: reguły zmiany i mapa wszystkich stanów}
\label{ch:motion}
\index{ruch}\index{równanie różniczkowe}

\noindent
Herbata z rozdziału~\ref{ch:models} żyła w pełnych minutach: jej reguła
przeskakiwała od jednej minuty do następnej. Wahadło nie przeskakuje; w każdej
chwili gdzieś jest i dokądś zmierza. Prawa ruchu Newtona, opublikowane w 1687
roku, nie mówią, gdzie będzie planeta w przyszły wtorek. Mówią, jak szybko
zmienia się jej prędkość w tej chwili. Jak wydobyć przyszłość z takiej
reguły?

Ten rozdział odpowiada sztuczką z rozdziału~\ref{ch:iteration}: rób małe
kroki. Potem rysuje obraz, w którym narysowana jest reszta książki \dash{}
mapę, na której każdy możliwy stan jest punktem, a każda przyszłość drogą
\dash{} i kończy się twierdzeniem, które mówi, dlaczego model pogody z
rozdziału~\ref{ch:lorenz} potrzebuje trzech liczb.

\begin{outcomes}
\outcome{czytać regułę zmiany jako informację o tym, jak szybko zmienia się w
  tej chwili każda liczba stanu}
\outcome{zamienić regułę zmiany na drogę w kodzie, metodą Eulera i metodą
  RK4}
\outcome{użyć energii, żeby sprawdzić, czy symulacja mówi prawdę}
\outcome{czytać portret fazowy: jego punkty stałe, jego pętle i kierunek
  ruchu}
\outcome{powiedzieć, co ruch na płaszczyźnie może, a czego nie może, i ilu
  liczb potrzebuje chaos w czasie ciągłym}
\end{outcomes}

\section{Reguła, która mówi, jak szybko}
\label{sec:ch04-rates}
\index{reguła zmiany}\index{prędkość}

\noindent
Prędkościomierz w samochodzie to reguła zmiany. Nie mówi, gdzie jesteś;
mówi, jak szybko zmienia się to, gdzie jesteś.

\begin{think}
Prędkościomierz pokazuje $60$ kilometrów na godzinę. Mniej więcej gdzie
będzie samochód za minutę? A za godzinę?
\end{think}
\reveal{Za minutę mniej więcej kilometr dalej: $60$ kilometrów na godzinę to
kilometr na minutę. Za godzinę nie da się tego powiedzieć, bo prędkość
wiele razy się zmieni.}
W krótkim czasie to, jak szybko coś się zmienia, mówi, jak daleko zajdzie; w
długim mówi niewiele. Więc przejeżdżasz kawałek, znowu patrzysz na
prędkościomierz i jedziesz dalej.

Oto pierwszy przykład. Ciężarek wisi na sprężynie; odciągasz go o jeden metr
od położenia spoczynku i puszczasz. Jego stan to dwie liczby:
\emph{położenie} $x$, czyli jak daleko jest od spoczynku, i \emph{prędkość}
$v$, czyli jak szybko zmienia się $x$. Reguła mówi, jak szybko zmienia się
każda z nich:
\begin{itemize}
\item położenie zmienia się w tempie $v$, bo tym właśnie jest prędkość;
\item prędkość zmienia się w tempie $-k\,x$: sprężyna ciągnie z powrotem ku
  środkowi, tym mocniej, im bardziej jest rozciągnięta. Liczba $k$ mówi, jak
  sztywna jest sprężyna; przyjmij $k = 1$. Znak minus mówi, że siła ciągnie
  wstecz.
\end{itemize}
Inne książki zapisują tę samą regułę jako
\[ \frac{\mathrm{d}x}{\mathrm{d}t} = v, \qquad
   \frac{\mathrm{d}v}{\mathrm{d}t} = -k\,x, \]
gdzie $\tfrac{\mathrm{d}x}{\mathrm{d}t}$ czyta się \enquote{jak szybko
zmienia się $x$}. Reguła tego rodzaju to \emph{równanie różniczkowe}. Ta
książka nigdy nie rozwiązuje go ręcznie. Śledzi je małymi krokami, tak jak
robi to komputer prognozujący pogodę.

\section{Małe kroki}
\label{sec:ch04-euler}
\index{metoda Eulera}\index{krok czasowy}

\noindent
Najprostszy sposób śledzenia reguły zmiany nosi imię Leonharda Eulera, który
opisał go w XVIII wieku. Wybierz krótki odcinek czasu, \emph{krok}, i przesuń
każdą liczbę o jeden krok w tempie, jakie ma teraz:
\[ \text{nowe } x = x + \text{krok} \times v, \qquad
   \text{nowe } v = v + \text{krok} \times (-k\,x). \]
Potem zapytaj regułę znowu w nowym stanie i powtarzaj. Reguła zmiany stała
się regułą następnego kroku i dalej działa iteracja.

\pyregion{code/ch04/spring.py}{euler}{Sprężyna śledzona metodą Eulera}{lst:ch04-euler}

\noindent
\code{rates} to reguła, a \code{euler\_step} to metoda Eulera, dwa wiersze.
\code{energy} dodaje energię ruchu, połowę $v \times v$, do energii
zmagazynowanej w sprężynie, połowy $k \times x \times x$. Ciężarek startuje z
$x = 1$, w spoczynku, a kroki trwają \val{ch04.euler.dt} sekundy.

\mermaidfig{ch04-step-loop}{Metoda Eulera. Reguła podaje tempa w tej chwili;
mały krok przesuwa tak, jakby się nie zmieniały; w nowym stanie znowu pyta
się regułę.}{reguła zmiany śledzona krokami}

\begin{think}
Prawdziwa sprężyna bez tarcia drga w nieskończoność, za każdym razem
dochodząc tak samo daleko: nic nie odbiera jej energii i nic jej nie dodaje.
Co powinna pokazywać kolumna energii, wiersz po wierszu?
\end{think}
\reveal{Za każdym razem tę samą liczbę: \num{0.5}, połowę $1 \times 1$.}

\transcript{ch04-spring-euler}

\noindent
Nie pokazuje. Wychylenia rosną, a energia wspina się od \num{0.5} do
\val{ch04.euler.energy.ten} w dziesięć sekund. Nic w regule nie dodało tej
energii. Zrobiła to symulacja.

Powód widać na obrazku. Narysuj stan jako punkt: położenie w poziomie,
prędkość w pionie. Przy $k = 1$ energia to połowa kwadratu odległości punktu
od środka, więc prawdziwa sprężyna krąży po okręgu. Krok Eulera idzie wzdłuż
prostej, którą w tej chwili wskazuje reguła, a prosta styczna do okręgu
wychodzi z niego na zewnątrz. Każdy krok ląduje trochę na zewnątrz, i to
dokładnie, z twierdzenia Pitagorasa: każdy krok mnoży energię przez
$1 + \text{krok} \times \text{krok}$, czyli \val{ch04.euler.factor} przy
kroku \val{ch04.euler.dt}. Jeden procent na krok, \val{ch04.euler.steps}
razy w dziesięć sekund.

\begin{think}
Weź kroki dziesięć razy mniejsze, \val{ch04.small.dt} sekundy. Czy symulacja
jest teraz poprawna?
\end{think}
\reveal{Poprawna dłużej, a nie poprawna. W dziesięć sekund energia rośnie o
\val{ch04.small.gain.ten} procent zamiast o \val{ch04.euler.gain.ten}; po
stu sekundach urosła o \val{ch04.small.gain.hundred} procent, mniej więcej
tyle, ile większy krok zyskał w dziesięć.}
Każdy mały krok zyskuje sto razy mniej, a w każdej sekundzie jest ich
dziesięć razy więcej. Dziesięć razy mniejsze kroki kupiły dziesięć razy
dłuższy czas, za dziesięć razy więcej pracy.

\begin{trapbox}
\textbf{\enquote{Mniejszy krok sprawia, że symulacja jest poprawna.}} Sprawia,
że jest poprawna dłużej. Błąd metody krokowej powstaje na nowo przy każdym
kroku i się piętrzy, więc mniejszy krok kupuje czas, zanim ta sterta urośnie
za bardzo. Uczciwe pytanie do symulacji nie brzmi, jak mały miała krok, tylko
jak długo ją sprawdzano. Część~II spotyka surowszy horyzont, którego nie
przesunie żaden rozmiar kroku.
\end{trapbox}

\section{Spójrz naprzód i uśrednij}
\label{sec:ch04-rk4}
\index{metoda Rungego--Kutty}\index{RK4}

\noindent
Euler pyta regułę raz, na początku kroku, i wierzy odpowiedzi przez cały
krok. Tempa zmieniają się w trakcie kroku, a Euler nigdy się o tym nie
dowiaduje. Metoda, której używa ta książka, opracowana przez Carla Rungego i
Martina Kuttę na przełomie XIX i XX wieku i znana jako RK4, pyta cztery razy:
\begin{enumerate}
\item o tempa na początku;
\item o tempa pół kroku dalej, idąc w pierwszych tempach;
\item o tempa pół kroku dalej, idąc w drugich tempach;
\item o tempa cały krok dalej, idąc w trzecich.
\end{enumerate}
Potem robi prawdziwy krok ze średnią tych czterech, licząc dwie środkowe
dwukrotnie.

\pyregion{code/src/chaoslab/flows.py}{rk4}{Przepis RK4 z biblioteki książki}{lst:ch04-rk4}

\noindent
\code{field} to reguła zmiany, \code{k1} do \code{k4} to cztery próbki, a
\code{\_add} wylicza, gdzie byłbyś po przejściu kawałka w danych tempach.
Wagi $1$, $2$, $2$, $1$ są dobrane tak, żeby błędy próbek znosiły się
nawzajem, jak tylko się da; algebrę tu pomijamy. Ta sama funkcja, obok
\api{euler\_step} i \api{integrate}, poniesie model pogody z
rozdziału~\ref{ch:lorenz}.

\begin{lab}{l04_01_rk4}{Jeden krok RK4}
Napisz \code{rk4\_step} według przepisu. Test porównuje go z biblioteką na
sprężynie i na wahadle, a także na regule \enquote{$x$ rośnie w tempie
równym $x$}: jeden krok długości $1$ z $1$ musi wylądować na $1 + 1 +
\tfrac{1}{2} + \tfrac{1}{6} + \tfrac{1}{24}$, blisko prawdziwej odpowiedzi
$e \approx \num{2.718}$.
\end{lab}

\begin{think}
RK4 pyta regułę cztery razy na krok, więc za tę samą pracę Euler może robić
kroki cztery razy mniejsze. Co wolisz: Eulera z krokami \val{ch04.same.dt}
sekundy czy RK4 z krokami \val{ch04.euler.dt}?
\end{think}
\reveal{RK4, z ogromną przewagą. W dziesięć sekund Euler zyskuje
\val{ch04.same.gain} procent energii; RK4 traci \val{ch04.rk4.ppm} części na
milion. Za tę samą pracę jego błąd jest około \val{ch04.rk4.better} razy
mniejszy.}

\pyregion{code/ch04/energy_drift.py}{drift}{Pomiar dryfu energii symulowanej sprężyny}{lst:ch04-drift}

\transcript{ch04-energy-drift}

\noindent
\code{rule calls} liczy pracę. Ostatni wiersz mówi, że RK4 też nie jest
dokładne: w tysiąc sekund traci \val{ch04.rk4.ppm.thousand} części na
milion. Ono również jest poprawne dłużej, a nie poprawne, tylko o wiele
dłużej.

\begin{lab}{l04_02_drift}{Zmierz dryf}
Napisz energię sprężyny i funkcję, która śledzi sprężynę daną metodą przez
dany czas i podaje, o ile zmieniła się jej energia. Potem odpowiedz na
pytanie, które rozdział zostawia otwarte: zmniejsz krok o połowę \dash{} ile
razy mniejszy stanie się dryf energii w każdej z metod? (Komentarz w listingu
RK4 dotyczy błędu położenia, a to inna miara.) Przewiduj, zanim uruchomisz.
\end{lab}

\plotfig{ch04-euler-rk4}{Ta sama sprężyna śledzona przez dziesięć sekund
dwiema metodami z tym samym krokiem. Po lewej: na płaszczyźnie stanów Euler
(czerwony) oddala się spiralą od prawdziwego okręgu, a RK4 (niebieski)
zostaje na nim. Po prawej: energia, którą podaje każda z
nich.}{Euler uciekający spiralą i RK4}

\section{Mapa wszystkich stanów}
\label{sec:ch04-phase}
\index{przestrzeń fazowa}\index{portret fazowy}\index{wahadło}

\noindent
Lewa połowa rysunku~\ref{fig:ch04-euler-rk4} nie pokazuje niczego w
funkcji czasu. Każdy jej punkt to cały stan, położenie i prędkość naraz, a
ruch to droga od punktu do punktu. Płaszczyzna wszystkich możliwych stanów to
\emph{przestrzeń fazowa}, a rysunek dróg na niej to \emph{portret fazowy}.
Każdy układ chaotyczny w tej książce jest narysowany na takim portrecie.

Bogatszym przykładem jest wahadło. Jego stan to kąt od pionu w dół i to, jak
szybko ten kąt się zmienia. Jego reguła to reguła sprężyny z jedną różnicą:
siła ciągnąca z powrotem jest największa, gdy wahadło sterczy w bok, a znika,
gdy wskazuje prosto w dół albo prosto w górę. (To sinus, którego woła reguła
wahadła z \code{chaoslab}.)

\begin{think}
Wahadło puszczone dalej ma dłuższą drogę, ale też porusza się szybciej. Czy
wahadło puszczone z $90$ stopni wychyla się tam i z powrotem dłużej niż
puszczone z $10$?
\end{think}
\reveal{Dłużej: \val{ch04.pend.ninety} sekundy wobec \val{ch04.pend.ten},
o około \val{ch04.pend.ninety.pct} procent więcej. Puszczone ze $170$
stopni potrzebuje \val{ch04.pend.big} sekundy, \val{ch04.pend.ratio} razy
więcej.}

\pyregion{code/ch04/pendulum.py}{period}{Mierzenie czasu wahnięcia wahadła}{lst:ch04-period}

\transcript{ch04-pendulum}

\noindent
Szkolna reguła, że czas wahnięcia nie zależy od wychylenia, jest prawdziwa
dla małych wychyleń: czasy zbliżają się do \val{ch04.pend.small} sekundy, gdy
wychylenie maleje. Dlatego precyzyjny zegar wahadłowy wychyla się ledwie o
kilka stopni.

Teraz narysuj regułę jako mapę. W każdym punkcie płaszczyzny reguła podaje
tempa, a więc kierunek, w którym stan zaraz się ruszy. Postaw tam małą
strzałkę. Ruch to podążanie za strzałkami.

\begin{think}
Gdzieś na mapie strzałka ma zerową długość: nic się nie zmienia. Jakie to
stany, jako położenia prawdziwego wahadła? Co się stanie, gdy każdy z nich
lekko trącisz?
\end{think}
\reveal{Zwisanie prosto w dół w spoczynku i balansowanie dokładnie do góry
nogami w spoczynku. Trącone, pierwsze łagodnie się wokół niego waha; drugie
odpada. To dwa punkty stałe, jeden stabilny, drugi niestabilny, jak w
rozdziale~\ref{ch:iteration}.}

\plotfig{ch04-pendulum-map}{Mapa wszystkich stanów wahadła o długości
jednego metra. Szare strzałki pokazują, w którą stronę rusza się stan.
Niebieskie pętle to wahania puszczone z 30, 90 i 150 stopni; zielone drogi
przechodzą przez górę; linia przerywana oddziela jedne od drugich. Pełna
kropka to zwisanie w spoczynku; dwa puste kółka to to samo położenie do góry
nogami.}{portret fazowy wahadła}

\noindent
Blisko środka drogi są zamkniętymi pętlami: wahadło wychyla się tam i z
powrotem, w nieskończoność po tej samej pętli. Droga, która zamyka się sama w
sobie, to \emph{cykl}, a cykl jest doskonale przewidywalny: jedno okrążenie i
znasz wszystkie następne. Daleko u góry i u dołu drogi biegną przez całą
szerokość: wahadło przechodzi przez górę i obraca się dalej.

\section{Zatrzymać się, krążyć albo uciec}
\label{sec:ch04-settle}
\index{atraktor}\index{cykl graniczny}\index{twierdzenie Poincarégo--Bendixsona}

\noindent
Prawdziwe wahadło ma tarcie, które z każdym wahnięciem odbiera trochę
energii.

\begin{think}
Co dzieje się na mapie z drogą wahadła z tarciem puszczonego ze $150$
stopni i gdzie się ona kończy?
\end{think}
\reveal{Każde wahnięcie jest mniejsze od poprzedniego, więc droga nie może się
zamknąć. Zwija się spiralą do środka i kończy na pełnej kropce: wahadło zwisa
prosto w dół, w spoczynku.}
Stan, do którego wpadają wszystkie pobliskie drogi, to \emph{atraktor}, a
punkt spoczynku wahadła z tarciem to jego najprostszy rodzaj. Przyciągać może
też pętla. Holenderski inżynier Balthasar van der Pol zapisał klasyczny
przykład mniej więcej sto lat temu, dla obwodów z lampami radiowymi:
\[ x \text{ zmienia się w tempie } v, \qquad
   v \text{ zmienia się w tempie } -x + (1 - x \times x)\, v. \]
Dopóki $x$ leży między $-1$ a $1$, nowy człon przyspiesza ruch; poza tym
przedziałem hamuje. Zacznij blisko środka, od \val{ch04.vdp.start.in}, a
wychylenia rosną; zacznij daleko, od \val{ch04.vdp.start.out}, a maleją. Po
\val{ch04.vdp.time} sekundach oba wychylają się do \val{ch04.vdp.reach.in} i
\val{ch04.vdp.reach.out}: jedna pętla, osiągnięta od środka i od zewnątrz.
To \emph{cykl graniczny}, a zegar wahadłowy jest jednym z nich: małe
pchnięcie przy każdym wahnięciu, tarcie odbierające energię i wychylenie,
które ustala się na jednej wielkości, jakkolwiek je zacząłeś.

\plotfig{ch04-settle}{Dwa sposoby, w jakie ruch na płaszczyźnie może się
ustalić. Po lewej: wahadło z tarciem puszczone ze 150 stopni zwija się
spiralą do punktu, w którym zwisa w spoczynku. Po prawej: reguła van der Pola
ze startu wewnątrz i ze startu na zewnątrz, nawijająca się na tę samą
pętlę.}{ustalanie się w punkcie i na pętli}

\begin{think}
Dlaczego dwie różne drogi na takiej mapie nigdy się nie przecinają?
\end{think}
\reveal{W punkcie przecięcia stan miałby dwie przyszłości, a reguła daje
jedną strzałkę: jeden stan, jedna przyszłość.}
Na płaszczyźnie ten jeden fakt ma wielką siłę. Droga, która nie może
przeciąć ani siebie, ani żadnej innej drogi, i nie opuszcza ograniczonego
obszaru, jest ogrodzona własną przeszłością. Może zwijać się do punktu
stałego albo do pętli, albo uciec do nieskończoności. Nie może błąkać się bez
końca.

\begin{rigourbox}
To \emph{twierdzenie Poincarégo--Bendixsona}, od nazwisk Henriego Poincarégo
i Szweda Ivara Bendixsona. Wymaga gładkiej reguły, a jego
dokładne sformułowanie dopuszcza jeszcze jedno zakończenie, łańcuch punktów
stałych połączonych drogami. Dowód, na kilka stron, jest w każdym podręczniku
układów dynamicznych. Dotyczy ruchu w czasie
ciągłym: reguła skacząca krokami może przeskoczyć własną drogę, a
rozdział~\ref{ch:logistic} znajduje dzikie zachowanie w regule krokowej z
jedną liczbą. I wymaga płaszczyzny: w trzech wymiarach droga może przejść
nad sobą i pod sobą.
\end{rigourbox}

\noindent
Chaos w czasie ciągłym potrzebuje więc co najmniej trzech liczb stanu. Dlatego
model pogody z rozdziału~\ref{ch:lorenz}, który Edward Lorenz okroił w 1963
roku z większego modelu konwekcji, ma trzy liczby, a nie dwie.

\begin{think}
Zawieś drugie wahadło u dołu pierwszego. Ilu liczb potrzebuje teraz stan i
czy twierdzenie nadal obiecuje prosty ruch?
\end{think}
\reveal{Czterech: kąta i prędkości dla każdego wahadła. Twierdzenie nic nie
mówi o czterech liczbach, a rozdział~\ref{ch:mechanics} pokazuje, że ruch
wahadła podwójnego jest chaotyczny.}

\begin{trapbox}
\textbf{\enquote{Wahadło to prosty układ, więc jego ruch musi być prosty.}}
Dla jednego wahadła tak jest: dwie liczby stanu, płaszczyzna i twierdzenie,
które pozwala tylko zatrzymać się, krążyć albo uciec. Zawieś na nim drugie, a
stan ma cztery liczby, płot znika i ruch może stać się chaotyczny. To, jak
prosto wygląda układ, mówi niewiele o tym, jak prosty jest jego ruch. Więcej
mówi policzenie liczb w jego stanie.
\end{trapbox}

\section{Dwa ciała: mechanizm zegara w najlepszym wydaniu}
\label{sec:ch04-planets}
\index{zagadnienie dwóch ciał}\index{trzecie prawo Keplera}

\noindent
Cztery liczby stanu umożliwiają chaos; nie sprawiają, że on następuje.
Najlepszym dowodem jest układ, na którym zbudowano obraz świata jak zegara:
jedna planeta krążąca wokół Słońca. Grawitacja Newtona to reguła zmiany.
Położenie planety zmienia się w tempie jej prędkości, a prędkość zmienia się
w stronę Słońca o wielkość, która maleje z kwadratem odległości: dwa razy
dalej, cztery razy słabiej. Słońce jest tak dużo cięższe, że może stać w
miejscu.

\pyregion{code/ch04/two_bodies.py}{gravity}{Grawitacja Newtona jako reguła zmiany}{lst:ch04-gravity}

\noindent
Stan to cztery liczby: gdzie jest planeta, w poziomie i w pionie, i jak
szybko się porusza, w poziomie i w pionie. Listing wystrzeliwuje ją w bok z
\val{ch04.kepler.orbits} różnymi prędkościami, śledzi każdą orbitę dziesięć
razy dookoła metodą RK4 i mierzy czas każdego powrotu na linię startu.
\code{size} orbity to połowa jej największej szerokości.

\transcript{ch04-two-bodies}

\noindent
Średnia z dziesięciu okrążeń to pierwsze okrążenie, co do każdej wydrukowanej
cyfry: planeta wraca po swojej elipsie zgodnie z rozkładem, raz za razem.
Ostatnia kolumna jest taka sama dla wszystkich trzech orbit,
\val{ch04.kepler.law}: kwadrat czasu okrążenia podzielony przez sześcian
rozmiaru. Johannes Kepler znalazł tę regułę w położeniach planet w 1619 roku,
a Newton pokazał w 1687, że wynika ona z grawitacji; twój laptop znalazł ją
ponownie z samej reguły zmiany.

Dlaczego przy czterech liczbach nie ma chaosu? Dwie wielkości nigdy się nie
zmieniają: energia i tempo, w jakim odcinek od Słońca do planety zakreśla
pole. Każda z nich mocniej przyszpila stan, a przy dwóch przyszpilonych ruch
nie ma już miejsca, żeby być czymkolwiek poza regularnym. Dodaj trzecie
ciało, a to przestaje być prawdą; od tego zaczyna się
rozdział~\ref{ch:mechanics}.

\begin{worldbox}
Tablice położeń planet, dużych księżyców i wielu planetoid, według których
sterują nawigatorzy sond kosmicznych, publikuje Jet Propulsion Laboratory
NASA. Powstają tak, jak ten listing wyznaczył swoje orbity: Newtonowska
reguła zmiany dla całego Układu Słonecznego śledzona naprzód małymi krokami i
sprawdzana z pomiarami radarowymi i śledzeniem sond.
\end{worldbox}

\noindent
To mechanizm zegara w najlepszym wydaniu: reguła zmiany, małe kroki i
przewidywanie, które wraca zgodnie z rozkładem co do każdej wydrukowanej
cyfry. Część~II pyta, co się dzieje, gdy reguła jest równie dokładna, a
przyszłość nie.

\begin{summarybox}
\begin{itemize}
\item \textbf{Reguła zmiany}, czyli \textbf{równanie różniczkowe}, mówi, jak
  szybko zmienia się w tej chwili każda liczba stanu, a nie jaka będzie za
  chwilę.
\item \textbf{Metoda Eulera} przesuwa każdą liczbę o jeden krótki krok w
  tempie, jakie ma teraz. Na sprężynie zyskuje energię przy każdym kroku, a
  mniejszy krok czyni ją \textbf{poprawną dłużej}, a nie poprawną.
\item \textbf{RK4} próbkuje tempa cztery razy na krok i je uśrednia; za tę
  samą pracę jej błąd jest dużo mniejszy, ale ona też dryfuje. Sprawdzianem
  jest energia.
\item \textbf{Przestrzeń fazowa} to płaszczyzna wszystkich możliwych stanów;
  reguła to strzałka w każdym punkcie. Wahadło rysuje \textbf{cykle} wokół
  swojego stabilnego \textbf{punktu stałego}, a duże wahnięcia trwają dłużej
  niż małe.
\item Tarcie zwija drogę spiralą do \textbf{atraktora}; przyciągać może też
  pętla. Drogi się nie przecinają, więc na płaszczyźnie ruch może się tylko
  zatrzymać, krążyć albo uciec (\textbf{Poincaré--Bendixson}): chaos w czasie
  ciągłym potrzebuje co najmniej trzech liczb.
\item Dwa ciała mają cztery liczby stanu, a mimo to kreślą elipsy zgodnie z
  rozkładem, jak mówi trzecie prawo Keplera: trzy liczby są potrzebne do
  chaosu, ale nie wystarczają.
\end{itemize}
\end{summarybox}

\canyou

\begin{checkpoint}
\item Prędkościomierz pokazuje $90$ kilometrów na godzinę. Mniej więcej jak
  daleko samochód zajedzie w ciągu następnych $20$ sekund?
  \answerto{Około pół kilometra: $90$ kilometrów na godzinę to \num{1.5}
  kilometra na minutę, a $20$ sekund to jedna trzecia minuty.}
\item Sprężyna z $k = 1$ jest w $x = 0$ i porusza się z $v = 1$, więc jej
  energia wynosi \num{0.5}. Zrób jeden krok Eulera o długości \num{0.1}.
  Jakie są nowe $x$, $v$ i energia?
  \answerto{$x = \num{0.1}$ i $v = 1$, więc energia wynosi
  $\num{0.5} + \num{0.005} = \num{0.505}$: o jeden procent więcej, czynnik
  $1 + \num{0.1} \times \num{0.1}$.}
\item Euler z krokiem \num{0.1} zyskuje w dziesięć sekund
  \val{ch04.euler.gain.ten} procent energii sprężyny. Mniej więcej jak
  długo może trwać przebieg z krokiem \num{0.001}, zanim zyska tyle samo?
  \answerto{Około tysiąca sekund: krok sto razy mniejszy zyskuje sto razy
  mniej na sekundę.}
\item Opisz cztery próbki, które RK4 pobiera w jednym kroku, i jak je łączy.
  \answerto{Tempa na początku; pół kroku dalej w pierwszych tempach; pół
  kroku dalej w drugich; cały krok dalej w trzecich. Przesuwa się ze średnią
  o wagach $1$, $2$, $2$, $1$.}
\item Co na portrecie fazowym wahadła oznacza dla prawdziwego wahadła
  zamknięta pętla, droga przez całą szerokość i spirala?
  \answerto{Wahanie się tam i z powrotem; przechodzenie przez górę i
  obracanie się; utratę energii przez tarcie i zatrzymanie się w zwisie.}
\item Dlaczego gładka reguła zmiany z dwiema liczbami stanu nigdy nie jest
  chaotyczna i ilu liczb potrzebuje chaos w czasie ciągłym?
  \answerto{Jej drogi nigdy się nie przecinają, bo każdy stan ma jedną
  przyszłość, więc ograniczona droga może tylko zatrzymać się w punkcie
  stałym albo nawinąć na pętlę (Poincaré--Bendixson). Chaos w czasie
  ciągłym potrzebuje co najmniej trzech.}
\end{checkpoint}
